\section{Analysis}
\label{sec:analysis}

We analyse the \NRecords{} responses to the \NQuestions{} questions answered in all 12 conditions (Section~\ref{sec:progress}). All differences are in percentage points. We can grade a response only if it ends with the requested \texttt{Answer:} line. \NCutoff{} responses (\NCutoffPct{}\%) have no such line because they reached the 4,096-token output limit and were cut off mid-reasoning; \NCutoffEnCot{} of them are English CoT. We leave these out of accuracy rather than count them as wrong, and check in Section~\ref{sec:cot-delta} that counting them as wrong changes no conclusion.

Differences between paired conditions are tested with the exact McNemar test \citep{mcnemar-1947-note} with Holm correction \citep{holm-1979-simple}, and confidence intervals (CIs) come from a 2,000-sample bootstrap that resamples questions, so that the 12 responses to a question stay together \citep{berg-kirkpatrick-etal-2012-empirical,dror-etal-2018-hitchhikers}. Square brackets give 95\% CIs, the range of values consistent with the data; an interval that excludes zero means the difference is unlikely to be due to chance.

\subsection{Does CoT help, and in which language?}
\label{sec:cot-delta}

\begin{table*}[t]
\centering
\begin{tabular}{lrrrrrrr}
\toprule
 & & \multicolumn{3}{c}{English} & \multicolumn{3}{c}{Hindi} \\
\cmidrule(lr){3-5}\cmidrule(lr){6-8}
Task & $Q$ & Direct & CoT & $\Delta$ & Direct & CoT & $\Delta$ \\
\midrule
Date addition & 600 & 93.7 & 98.5 & +4.8$^{*}$ & 89.5 & 88.3 & $-$1.2\phantom{$^{*}$} \\
Date subtraction & 450 & 93.6 & 95.9 & +2.3$^{*}$ & 92.2 & 78.2 & \textbf{$-$14.0$^{*}$} \\
Date duration & 585 & 88.9 & 99.1 & +10.2$^{*}$ & 85.1 & 87.5 & +2.5\phantom{$^{*}$} \\
Date recurrence & 450 & 24.5 & 92.4 & +67.9$^{*}$ & 4.8 & 26.0 & +21.2$^{*}$ \\
Day of week & 445 & 14.1 & 70.5 & +56.4$^{*}$ & 15.7 & 21.8 & +6.1$^{*}$ \\
Interval date & 150 & 29.8 & 43.1 & +13.3$^{*}$ & 35.1 & 17.9 & \textbf{$-$17.2$^{*}$} \\
Time addition & 149 & 57.7 & 97.1 & +39.3$^{*}$ & 60.2 & 73.5 & +13.4$^{*}$ \\
Time subtraction & 111 & 48.9 & 98.5 & +49.5$^{*}$ & 51.7 & 74.2 & +22.5$^{*}$ \\
Time duration & 149 & 58.8 & 99.8 & +40.9$^{*}$ & 28.9 & 72.6 & +43.8$^{*}$ \\
\midrule
All tasks & 3,089 & 63.2 & 91.1 & +27.9$^{*}$ & 57.7 & 62.9 & +5.1$^{*}$ \\
\bottomrule
\end{tabular}
\caption{Accuracy (\%) with direct answers and chain-of-thought (CoT), pooled over difficulty levels and distractor conditions. $Q$ is the number of questions; each question contributes three responses per language and condition. $\Delta$ is CoT minus direct in percentage points; $^{*}$ marks $p<0.05$ in an exact McNemar test on paired responses, Holm-corrected over the 18 task--language tests (uncorrected for the pooled row). Significant negative deltas are in bold. Chance for day of week is 14.3\%.}
\label{tab:cot-delta}
\end{table*}


Table~\ref{tab:cot-delta} shows the central result. CoT raises English accuracy from \EnDirectAcc{}\% to \EnCotAcc{}\% (\EnCotDelta{} points [\EnCotDeltaLo{}, \EnCotDeltaHi{}]) but Hindi accuracy only from \HiDirectAcc{}\% to \HiCotAcc{}\% (\HiCotDelta{} [\HiCotDeltaLo{}, \HiCotDeltaHi{}]). The English gain is significant in all \NEnSigPos{} tasks. In Hindi, CoT helps significantly in \NHiSigPos{} tasks and \emph{hurts} significantly in two: date subtraction (\HiSubtractionDirect{}\% $\rightarrow$ \HiSubtractionCot{}\%) and interval date (\HiIntervalDirect{}\% $\rightarrow$ \HiIntervalCot{}\%). Asking the model to reason step by step in Hindi can therefore be worse than asking it for the answer alone, and this happens without any distractor, since the table pools all three distractor conditions. The only task where Hindi gains at least as much as English is time duration (\HiTimeDurationDelta{} vs.\ \EnTimeDurationDelta{} points), where Hindi direct answers are unusually poor (\HiTimeDurationDirect{}\%).

These conclusions do not depend on how the cut-off responses are treated. Counting them as wrong, which penalises English CoT most, the CoT gain is still \EnStrictDelta{} points in English and \HiStrictDelta{} in Hindi.

\paragraph{Direct answers on two tasks are at chance.}
\input{figures/day_of_week_float}
Without CoT, day-of-week accuracy is \EnDowDirect{}\% in English and \HiDowDirect{}\% in Hindi, against a chance level of 14.3\%. Figure~\ref{fig:dow} shows that this is not a near miss. The difference between the predicted and the correct weekday is spread evenly over all seven values (\EnDowDirectMin{}--\HiDowDirectMax{}\% each), so the direct answer carries no information about the date. Instead, direct answers favour one weekday regardless of the question: \EnDowModeDay{} for \EnDowModeShare{}\% of English answers and \HiDowModeDay{} for \HiDowModeShare{}\% of Hindi answers, while the correct answers are spread roughly evenly over the week. English CoT fixes most of this (\EnDowCot{}\%), with remaining errors skewed late (+1 day: \EnDowCotPlusOne{}\%; $-$1 day: \EnDowCotMinusOne{}\%). Hindi CoT is more often exactly one day late (\HiDowCotPlusOne{}\%) than correct (\HiDowCotZero{}\%).

Date recurrence shows a similar split. In both languages, \EnRecDirectFirstStep{}\% of direct answers are ``today plus one interval'', which ignores which occurrence the question asks for. English CoT almost eliminates this shortcut (\EnRecCotFirstStep{}\%), while Hindi CoT keeps it in \HiRecCotFirstStep{}\% of answers. Hindi also produces a date in the past, today minus the interval times the occurrence number, in \HiRecDirectMirror{}\% (direct) and \HiRecCotMirror{}\% (CoT) of answers, against \EnRecDirectMirror{}\% and \EnRecCotMirror{}\% in English. The Hindi template phrases the question as \foreignlanguage{hindi}{आज को छोड़कर, 1वीं आवृत्ति की तिथि क्या है?} \emph{\={a}j ko cho\d{r}kar, 1v\={\i}\d{m} \={a}v\d{r}tti k\={\i} tithi ky\={a} hai?} (`excluding today, what is the date of the 1th occurrence?'), with a templated ordinal suffix that is non-standard for small numbers. We have not yet tested whether this wording causes the reversal; it is a possible confound in TRD's Hindi recurrence questions rather than in the model's arithmetic.

\subsection{The English--Hindi gap across difficulty}
\label{sec:difficulty}

\input{figures/difficulty_float}

We call English accuracy minus Hindi accuracy the \emph{English--Hindi gap}; a positive gap means English is ahead. Because each difficulty level contains a different mix of tasks, Figure~\ref{fig:difficulty} compares levels using only tasks present at all of them. From the short to the very long level, the gap grows far more with CoT than with direct answers. For date addition and duration it grows from \GapADirectShort{} to \GapADirectVeryLong{} points with direct answers but from \GapACotShort{} to \GapACotVeryLong{} points with CoT. Over the five date tasks, the CoT gap grows from \GapCotShort{} to \GapCotVeryLong{} points, while the direct gap stays below 9 points. On the easiest questions, Hindi CoT does not even beat Hindi direct answers (\AHiCotShort{}\% vs.\ \AHiDirectShort{}\%). The language gap is therefore mainly a gap in how much CoT helps.

\subsection{CoT under distraction}
\label{sec:distraction}

\begin{table}[t]
\centering
\small
\setlength{\tabcolsep}{4pt}
\begin{tabular}{lrrrrr}
\toprule
 & \multicolumn{3}{c}{Accuracy} & \multicolumn{2}{c}{Drop} \\
\cmidrule(lr){2-4}\cmidrule(lr){5-6}
 & None & Sim. & Dis. & Sim. & Dis. \\
\midrule
\multicolumn{6}{l}{\textit{English}} \\
\quad Direct & 63.9 & 63.0 & 62.7 & +0.9 & +1.2 \\
\quad CoT & 91.6 & 90.6 & 91.2 & +0.9 & +0.4 \\
\quad CoT $-$ direct &  &  &  & +0.0 & $-$0.8 \\
\midrule
\multicolumn{6}{l}{\textit{Hindi}} \\
\quad Direct & 57.4 & 58.1 & 57.7 & $-$0.7 & $-$0.3 \\
\quad CoT & 65.2 & 60.7 & 62.7 & +4.5 & +2.5 \\
\quad CoT $-$ direct &  &  &  & \textbf{+5.2} & \textbf{+2.8} \\
\bottomrule
\end{tabular}
\caption{Accuracy (\%) with no distractor (None), a time-related (Sim.) or an unrelated (Dis.) distractor. Drop is accuracy without the distractor minus accuracy with it (points; positive means the distractor hurt). CoT $-$ direct is the CoT drop minus the direct drop (the CoT distraction penalty); bold marks a 95\% cluster-bootstrap interval that excludes zero.}
\label{tab:distractor}
\end{table}


Our main hypothesis was that CoT gives a distractor more opportunities to enter the reasoning and change the answer. Table~\ref{tab:distractor} supports this only in Hindi. In English, neither distractor lowers accuracy by more than \EnMaxDrop{} points under either prompting condition. In Hindi, direct answers are unaffected by the similar distractor (accuracy changes by \HiDropDirectSim{} points), but CoT accuracy falls by \HiDropCotSimAbs{} points. CoT therefore loses \HiDidSimAbs{} more points to the distractor than direct answering does [\HiDidSimLo{}, \HiDidSimHi{}]; we call this extra loss the \emph{CoT distraction penalty} (the ``CoT $-$ direct'' row of Table~\ref{tab:distractor}). In English the penalty is \EnDidSim{} [\EnDidSimLo{}, \EnDidSimHi{}]. The unrelated distractor also carries a smaller penalty in Hindi (\HiDidDis{} [\HiDidDisLo{}, \HiDidDisHi{}]). Computing the Hindi penalty for the similar distractor separately for each task shows that it comes mostly from date duration (\HiDidSimDuration{} points [\HiDidSimDurationLo{}, \HiDidSimDurationHi{}]), day of week (\HiDidSimDow{} [\HiDidSimDowLo{}, \HiDidSimDowHi{}]) and time duration (\HiDidSimTimeDuration{} [\HiDidSimTimeDurationLo{}, \HiDidSimTimeDurationHi{}]). Previous work found that irrelevant context distracts models on arithmetic word problems \citep{shi-etal-2023-distracted}; here the susceptibility is specific to one language and to the reasoning condition.

\paragraph{Does the distractor enter the trace?}
We count a trace as mentioning the distractor if it contains a content word of the distractor that does not occur in the question (function words and generic time words such as ``day'' are excluded). Some of these words can appear in a trace by coincidence, for example month names in a trace about dates. To measure this chance rate we use a \emph{placebo}: for each question we take the trace written without any distractor and search it for the words of the distractor that would have been inserted. The model never saw those words, so any match there is coincidence. The difference between the two rates is the share of traces that really take up the distractor. Hindi traces take up the distractor far more often than English ones (Table~\ref{tab:trace}): \HiMentionDis{}\% of Hindi traces mention the unrelated distractor (placebo \HiMentionDisPlacebo{}\%), against \EnMentionDis{}\% in English (placebo \EnMentionDisPlacebo{}\%). Hindi traces often build the distractor into the narrative of the solution. For the short date-addition question that starts ``The garden needs new flowers'', a correct Hindi trace ends with
\begin{quote}
\foreignlanguage{hindi}{इस तरह, बगीचे में नए फूल लगाने के लिए सही तारीख 2025-03-30 है।}\\
\emph{is tarah, bag\={\i}ce me\d{m} nae ph\={u}l lag\={a}ne ke lie sah\={\i} t\={a}r\={\i}kh 2025-03-30 hai.}\\
`Thus, the correct date for planting new flowers in the garden is 2025-03-30.'
\end{quote}
Appendix~\ref{app:mention} relates these mentions to answer correctness. The distractor also leaks into wrong answers. For example, after the distractor ``February is such a unique month with its varying days'', a Hindi direct answer to a recurrence question whose correct answer is \mbox{2025-04-10} was \mbox{2025-02-06}. Overall, \HiReuseDirect{}\% (\HiReuseDirectN{}) of wrong Hindi direct answers contain a number, month or weekday taken from the similar distractor. With the same placebo as above (wrong answers to the same questions asked without the distractor, checked against the distractor that would have been inserted), this happens in only \HiReuseDirectPlacebo{}\% (\HiReuseDirectPlaceboN{}), so the difference is not chance (Fisher's exact test, \HiReuseDirectP{}). In English the two rates do not differ (\EnReuseDirect{}\% vs.\ \EnReuseDirectPlacebo{}\%, \EnReuseDirectP{}). In Hindi direct answers, the distractor therefore changes what the wrong answers are without changing how many there are.

\subsection{Properties of the reasoning traces}
\label{sec:traces}

\begin{table}[t]
\centering
\small
\setlength{\tabcolsep}{4pt}
\begin{tabular}{lrr}
\toprule
 & English & Hindi \\
\midrule
Cut off before answering & 3.0\% & 0.7\% \\
Median length (words) & 103 & 79 \\
Median completion tokens & 327 & 198 \\
Target-language ratio & 1.000 & 0.986 \\
Ratio below 0.9 & 0.0\% & 2.4\% \\
Has a verification step & 79.1\% & 1.5\% \\
Answer value in trace & 99.9\% & 99.0\% \\
Mentions sim.\ distractor & 15.1\% (10.4) & 31.8\% (8.3) \\
Mentions dis.\ distractor & 18.7\% (8.5) & 35.9\% (1.4) \\
\bottomrule
\end{tabular}
\caption{Properties of CoT responses. All rows except the first are computed over graded responses. The target-language ratio is the number of target-language words divided by the number of Hindi plus English words, with both counts summed over all responses before dividing, so longer traces weigh more. Distractor mention rates are for traces whose prompt contained that distractor; the value in parentheses is the placebo rate for the same questions without it.}
\label{tab:trace}
\end{table}


\paragraph{The model stays in Hindi, and drifting into English helps.}
The target-language ratio (Table~\ref{tab:trace}), the share of a trace's words that are in Hindi, is \HiRatioPooled{} for Hindi traces, and only \HiRatioBelow{}\% of them contain more than 10\% English words (ratio below 0.9). We call these \emph{drifting} traces. The Hindi results above therefore reflect reasoning written in Hindi. The \HiNRatioBelow{} drifting traces look less accurate overall (\HiAccRatioBelow{}\% vs.\ \HiAccRatioAbove{}\%), but mainly because of which questions they occur on: \DriftDow{} of them are day-of-week questions, one of the two hardest tasks for Hindi CoT. Compared within the same task, drifting traces are more accurate: \DriftDowAcc{}\% vs.\ \DriftDowAccNo{}\% on day of week and \DriftRecAcc{}\% vs.\ \DriftRecAccNo{}\% on date recurrence. Combining all tasks this way, a drifting trace has about twice the odds of a correct answer (Appendix~\ref{app:drift}). This agrees with the finding that English intermediate steps can match or beat steps in the question's language \citep{shi-etal-2023-mgsm}, and it motivates the cross-lingual prompting condition \citep{qin-etal-2023-cross} that we run next.

\paragraph{Hindi traces are shorter and rarely check their work.}
Hindi traces are shorter than English ones (median \HiMedianWords{} vs.\ \EnMedianWords{} words; \HiMedianTokens{} vs.\ \EnMedianTokens{} completion tokens). \EnVerification{}\% of English traces contain a verification step (``let me verify'', ``alternatively''), against \HiVerification{}\% of Hindi traces. Verification does not explain the English advantage, however. English traces with a verification step look less accurate overall (\VerAccWith{}\% vs.\ \VerAccWithout{}\%), but again because verification is most common in the hardest tasks; compared within the same task, verifying makes no clear difference to accuracy (Appendix~\ref{app:verification}).

\paragraph{Answers follow the trace.}
The final answer value appears in the trace in \EnAnswerInTrace{}\% of English and \HiAnswerInTrace{}\% of Hindi CoT responses, and only \EnWrongEndsGold{} of \EnWrong{} wrong English answers and \HiWrongEndsGold{} of \HiWrong{} wrong Hindi answers come from a trace whose last value was the correct one. Wrong answers are therefore wrong in the trace, not copied wrongly at the end. This is consistency between trace and answer, not evidence that the trace reflects the model's internal computation \citep{turpin-etal-2023-unfaithful}.

\subsection{Language-specific failures and error types}
\label{sec:errors}

Every question was asked in both languages, so we can pair the English and Hindi responses to the same question, distractor and prompting condition and ask whether the two languages fail on the same questions. Under CoT, of \AgreeCotPairs{} pairs, both are right in \AgreeCotBothRight{} and both wrong in \AgreeCotBothWrong{}; when they disagree, English is almost always the one that is right (\AgreeCotEnOnly{} pairs, against \AgreeCotHiOnly{} where only Hindi is right). If Hindi failed mainly on questions that are hard in any language, the two languages would tend to fail together. We measure this with the $\phi$ coefficient, a correlation between two right/wrong outcomes (0: unrelated, 1: always the same), computed within each task and difficulty so that differences between tasks do not inflate it. With direct answers the median $\phi$ is \AgreeDirectPhi{}: the languages partly fail on the same, harder questions. With CoT it is \AgreeCotPhi{}: whether Hindi CoT is right has almost nothing to do with whether English CoT is right on the same question. Hindi CoT failures are therefore specific to Hindi rather than a sign of hard questions.

The errors also differ in kind (Figure~\ref{fig:errors} in Appendix~\ref{app:errors}). \EnErrAllOffByOne{}\% of wrong English CoT answers are off by one day (\EnErrAddOffByOne{}\% in date addition and \EnErrSubOffByOne{}\% in date subtraction), the slip of an otherwise correct procedure. Only \HiErrAllOffByOne{}\% of wrong Hindi CoT answers are off by one. The rest include month slips, where the day and year are right but the month is wrong (\HiErrAddMonthSlip{}\% of Hindi date-addition errors, none in English), and the reversed direction in date recurrence discussed in Section~\ref{sec:cot-delta}. Hindi CoT does not just make more of the same mistakes as English CoT; it often follows a different, wrong procedure.
